\documentclass[12pt]{article}
\usepackage[utf8]{inputenc}
\usepackage[T1]{fontenc}
\usepackage{hyperref}
\title{Pye Krumen Language Resources for K-12 Curriculum Development}
\author{Compiled by Claude Sonnet 5.5 (claude-sonnet-5-5) for Andr\'as Kornai}
\date{}
\begin{document}
\maketitle

\section{Introduction}
Pye Krumen (ISO 639-3 pye; Glottocode pyek1235) is a Kru (Grebo group) language of southwestern C\^ote d'Ivoire, around Tai, B\'er\'eby and San Pedro. Sources describe it as a set of dialects (including Trepo, Wlepo, Yapo and Pie) with an estimated 20,000 local speakers. Little educational material was found.

\section{Primary Educational Resources}
No primer or textbook was found in the searches.

\section{Linguistic and Reference Materials}
\begin{itemize}
\item Webonary dictionary site for Pye (introduction page with classification, dialects and speaker estimates). \url{https://www.webonary.org/pye/overview/introduction/?lang=en}
\item Max Planck Institute numeral database entry for Pye Krumen, based on 2013--2014 fieldwork. \url{https://lingweb.eva.mpg.de/channumerals/Krumen-Pye.htm}
\end{itemize}

\section{Religious and Community Resources}
\begin{itemize}
\item Joshua Project profile of the Northeastern Krumen. \url{https://joshuaproject.net/people_groups/14003}
\item Global Recordings Network page for the Yapo variety. \url{https://www.globalrecordings.net/en/language/12363}
\item Global Recordings Network page for the Wlepo variety. \url{https://www.globalrecordings.net/en/language/12361}
\end{itemize}

\section{Digital and Audio Resources}
The Webonary site above is the main digital resource found. Global Recordings pages are listed above; the search summary indicated no recordings for the overall variety.

\section{Curriculum Development Context}
A dictionary exists online, which is a starting point for vocabulary work. No school materials or programmes were found; a curriculum would require new materials developed with the community and an agreed orthography.

\end{document}